\chapter{ANALYSE DES RESULTATS OBTENUS ET DISCUSSIONS}
\thispagestyle{empty}

\section{Pourquoi la revision du protocole importe}

L'enjeu interpretatif central du projet est que le protocole d'evaluation
corrige modifie le recit issu de la premiere ligne d'evaluation. Dans le
dispositif initial, Node2Vec non corrige semblait dominer tous les
concurrents, avec 16.06\% de Hit@10 et 0.0751 de MRR. Apres correction
methodologique, la famille la plus solide n'est plus la recuperation purement
structurelle, mais la famille hybride corrigee.

Ce basculement n'est pas une fluctuation mineure. Il montre que la conception
de l'evaluation peut modifier la conclusion d'une etude de recommandation sur
graphe. Le protocole corrige represente donc non seulement un nouveau cadre de
mesure, mais aussi une correction d'interpretation.

\section{Les signaux semantiques et structurels sont complementaires}

Les tableaux complets de resultats suggerent un schema coherent:
\begin{itemize}
    \item la recuperation semantique seule est utile et place souvent de bons
    candidats en tete de liste;
    \item la recuperation structurelle seule est elle aussi tres puissante,
    surtout pour la discrimination globale sur graphe;
    \item le comportement de recommandation le plus fort apparait lorsque les
    deux signaux sont combines avec soin.
\end{itemize}

C'est pourquoi la meilleure famille globale de recommandation dans le protocole
corrige n'est ni BGE-M3 pur ni Node2Vec pur, mais les hybrides
semantico-structurels corriges.

\section{L'evaluation corrigee}

\subsection{Qualite de recommandation en tete de liste}

Si l'objectif est d'optimiser ce qui apparait au sommet d'une liste classee
courte, les hybrides corriges sont les modeles les plus convaincants. Le duo
le plus solide en pratique est:
\begin{itemize}
    \item \textit{BGE pooled + Node2Vec (corrected)} comme meilleur modele sur
    Hit@10;
    \item \textit{BGE-M3 + Node2Vec (corrected)} comme modele legerement plus
    fort sur plusieurs metriques secondaires de classement.
\end{itemize}

Cette lecture est bien alignee avec l'objectif reel de l'application, car les
utilisateurs n'inspectent que la partie superieure de la sortie classee.

\subsection{Discrimination globale des liens}

Si l'objectif devient au contraire de separer les aretes positives et
negatives sur un grand nombre de paires scorees, le constat change. Sous
ROC-AUC et PR-AUC, les modeles purement structurels et leurs voisins legerement
fusionnes non corriges restent extremement forts. Node2Vec reste meme en tete
de la couche AUC.

Cela ne contredit pas l'analyse du sommet de liste. Cela signifie simplement
que la discrimination globale des liens et la qualite de recommandation ne sont
pas des objectifs identiques.

\section{Qualite de recommandation au-dela de la precision}

L'un des ajouts les plus utiles du cadre corrige est la nouvelle famille de
metriques de qualite de recommandation. Les rapports precedents insistaient
surtout sur la precision du classement, mais les analyses ulterieures montrent
que des modeles proches en Hit@10 peuvent tout de meme se comporter de maniere
tres differente:
\begin{itemize}
    \item GCN Only est faible du point de vue du classement, mais fort en
    diversite et en faible biais de popularite;
    \item BGE-M3 Only est plus fortement tire par la popularite et plus
    concentre;
    \item les hybrides corriges occupent le meilleur compromis entre precision
    et qualite de recommandation.
\end{itemize}

Cette perspective elargie est essentielle pour le projet. Un systeme de
recommandation qui renvoie sans cesse les noeuds les plus evidents ou les plus
populaires peut paraitre satisfaisant selon certaines metriques de classement,
mais il est moins utile comme outil d'exploration pour l'apprentissage.

\section{Pooling et correction de degre}

L'extension des methodes initales clarifie deux questions de conception supplementaires.

\subsection{Pooling}

La reduction de BGE-M3 a 128 dimensions avant concatenation a ete introduite
pour equilibrer la contribution dimensionnelle des vecteurs semantiques et
structurels. L'effet est reel mais limite. Le pooling aide legerement sur
Hit@10, sans pour autant dominer toutes les autres metriques secondaires.
L'interpretation correcte est donc que le pooling est utile, mais non
transformant.

\subsection{Correction de degre}

La correction de degre se comporte de facon tres differente selon la famille de
modeles. Dans le protocole corrige, elle n'a presque aucun effet sur Node2Vec
pur, mais elle compte de maniere significative pour les hybrides. C'est l'une
des lecons les plus claires de la phase tardive d'evaluation: une strategie de
normalisation ne doit pas etre jugee isolément, car son effet depend du contexte
de fusion dans lequel elle est utilisee.

\section{Evolution de la couche de parcours d'apprentissage}

La couche LLM a egalement change de facon notable apres l'etape des rapports
intermediaires. L'idee initiale etait d'utiliser un petit modele local via
Ollama pour une generation hierarchique de parcours entierement locale. En
pratique, cette approche s'est revelee trop fragile pour produire des sorties
structurees stables sur l'ensemble complet des recommandations top-20. Le
probleme ne concernait pas seulement la qualite factuelle, mais aussi la forme
de la sortie: JSON mal formes, affectations incompletes et classifications
semantiquement faibles etaient frequents.

La refonte ulterieure de l'application a conserve un moteur de recuperation
local, fonde sur les artefacts du projet mis en cache, mais a remplace
l'organisateur par un flux de travail a sortie structuree plus robuste utilisant
Gemini. Cela n'a pas transforme le LLM en modele central de classement. Au
contraire, l'organisateur a ete ramene a une couche de post-traitement
contrainte, dont la sortie est verifiee par des regles strictes cote
application. Il s'agit d'un usage plus realiste et plus defendable d'un LLM
dans ce projet. Dans cette architecture, le repli heuristique n'est pas
accessoire: il fait partie de la fiabilite du systeme, car il permet de
conserver une sortie presentable et ordonnee lorsque la generation structuree
ne respecte pas les contraintes imposees.

\section{Interpretation pratique}

D'un point de vue pratique, le projet conduit a une conclusion en deux volets:
\begin{itemize}
    \item si l'objectif est une discrimination large des aretes du graphe, les
    methodes purement structurelles restent extremement fortes;
    \item si l'objectif est la qualite de recommandation d'articles au sommet
    de la liste, les modeles hybrides semantico-structurels corriges
    constituent la meilleure direction globale.
\end{itemize}

Comme l'objectif reel du produit est la recommandation et non la seule
recuperation abstraite de liens, c'est la seconde interpretation qui doit
guider le choix final du systeme.
